% ================================================================
% SECTION 5: PROBLEM MODULE
% ================================================================
\section{Problem Module}
\label{sec:problem}

\begin{center}
\code{bionetflux.core.problem}
\end{center}

The \code{Problem} class (648 lines) is the base class for all 1D PDE problems.  Each domain in the network is assigned its own \code{Problem} instance.

% ------------------------------------------------------------------
\subsection{Constructor}

\begin{lstlisting}[language=Python, caption={Problem constructor}]
class Problem:
    def __init__(self, 
                 neq: int = 2,
                 domain_start: float = 0.0,
                 domain_length: float = 1.0,
                 parameters: np.ndarray = None,
                 problem_type: str = "keller_segel",
                 name: str = "unnamed_problem"):
\end{lstlisting}

If \code{parameters} is \code{None}, it defaults to \code{np.array([1.0, 1.0, 0.0, 0.0])}.  The attribute \code{domain\_end} is computed as \code{domain\_start + domain\_length}.

% ------------------------------------------------------------------
\subsection{Attributes}

\begin{longtable}{p{3.3cm}p{4.7cm}p{6.5cm}}
\toprule
\textbf{Attribute} & \textbf{Type} & \textbf{Description} \\
\midrule
\code{neq} & \code{int} & Number of equations \\
\code{domain\_start} & \code{float} & Parameter-space start \\
\code{domain\_length} & \code{float} & Parameter-space length \\
\code{domain\_end} & \code{float} & \code{domain\_start + domain\_length} \\
\code{name} & \code{str} & Human-readable name \\
\code{parameters} & \code{np.ndarray} & Physical parameter vector \\
\code{n\_parameters} & \code{int} & \code{len(parameters)} \\
\code{type} & \code{str} & Problem type string (e.g.\ \code{"keller\_segel"}, \code{"organ\_on\_chip"}) \\
\code{u\_names} & \code{List[str]} & Equation names; \code{['u','phi']} for \code{neq=2}, else \code{['u0','u1',...]} \\
\code{unknown\_names} & \code{List[str]} & Display names: \code{["Unknown n.\ 1", ...]} \\
\code{extrema} & \code{List[Tuple]} & Physical-space endpoints, default \code{[(domain\_start, 0), (domain\_end, 0)]} \\
\midrule
\multicolumn{3}{l}{\textit{Callable attributes (functions to be set by the user)}} \\
\midrule
\code{chi} & \code{Callable} & Chemotactic sensitivity $\chi(x)$ \\
\code{dchi} & \code{Callable} & Derivative $\chi'(x)$ \\
\code{force} & \code{List[Callable]} & Source terms $f_k(s,t)$, one per equation \\
\code{u0} & \code{List[Callable]} & Initial conditions $u_k^0(s)$ \\
\code{solution} & \code{List[Callable]} & Analytical solutions $u_k^{\mathrm{ex}}(s,t)$ \\
\code{flux\_solution} & \code{List[Optional[Callable]]} & Analytical flux solutions (may be \code{None}) \\
\code{flux\_u0} & \code{List[Callable]} & Left boundary fluxes $\sigma_k(t)$ at $s_0$ \\
\code{flux\_u1} & \code{List[Callable]} & Right boundary fluxes $\sigma_k(t)$ at $s_1$ \\
\code{neumann\_data} & \code{np.ndarray} & Shape $(4,)$ --- legacy Neumann data \\
\bottomrule
\end{longtable}

All callable lists are initialised with zero functions of the appropriate signature.

% ------------------------------------------------------------------
\subsection{Setter Methods}

\begin{longtable}{p{8cm}p{7cm}}
\toprule
\textbf{Method} & \textbf{Description} \\
\midrule
\code{set\_chemotaxis(chi, dchi)} & Set $\chi(x)$ and $\chi'(x)$ \\
\code{set\_force(equation\_idx, force\_func)} & Set source $f_k(s,t)$ \\
\code{set\_solution(equation\_idx, sol\_func)} & Set analytical solution $u_k^{\mathrm{ex}}(s,t)$ \\
\code{set\_flux\_solution(equation\_idx, flux\_func)} & Set analytical flux \\
\code{set\_initial\_condition(equation\_idx, u0\_func)} & Set $u_k^0$; auto-wraps 0-, 1-, or 2-argument functions to a unified signature \\
\code{set\_boundary\_flux(equation\_idx, left, right)} & Set left/right boundary fluxes \\
\code{set\_extrema(point1, point2)} & Set physical-space endpoint tuples \\
\code{get\_extrema() -> List[Tuple]} & Retrieve extrema \\
\code{set\_parameter(index, value)} & Set single parameter \\
\code{get\_parameter(index) -> float} & Get single parameter \\
\code{set\_parameters(params)} & Replace entire parameter vector \\
\code{set\_function(name, func)} & Dynamically set any callable attribute by name \\
\bottomrule
\end{longtable}

\begin{notebox}
\code{set\_initial\_condition} uses \code{inspect.signature} to determine the number of parameters of the supplied function and wraps it into a universal \code{f(s)} signature.  Functions with zero arguments are treated as constants.
\end{notebox}

% ------------------------------------------------------------------
\subsection{Validation and Testing}

\begin{longtable}{p{7cm}p{8cm}}
\toprule
\textbf{Method} & \textbf{Description} \\
\midrule
\code{validate\_problem(verbose=False) -> bool} & Checks: positive \code{neq}, positive \code{domain\_length}, parameter consistency, function list lengths, extrema format, problem type. Returns \code{True} if no errors.  The recognised problem types are \code{"keller\_segel"}, \code{"organ\_on\_chip"}, and \code{"generic"}. \\
\code{test\_functions(verbose=False) -> bool} & Calls every callable attribute on test arrays and checks return types/shapes. \\
\code{run\_self\_test(verbose=True) -> bool} & Runs validation, function tests, parameter get/set, extrema get/set, and dynamic function setting. \\
\code{create\_test\_problems() -> Dict[str, Problem]} & Class method returning a dictionary of test problems: \code{"ks\_basic"}, \code{"ooc\_basic"}, \code{"custom\_geometry"}, \code{"analytical"}, \code{"invalid"}. \\
\bottomrule
\end{longtable}
